\chapter*{Introduction}
\label{ch:introduction}
\addcontentsline{toc}{chapter}{\nameref{ch:introduction}}
\chaptermark{Introduction}
\section*{Affine manifolds}
If one forgets to bring their ruler to Euclidean space, they may stumble
upon the notion of an \emph{affine geometry}. Although a ruler has been lost,
an affine geometry still maintains the notion of straight lines and constant
speed. An \emph{affine manifold} is a manifold that is locally modelled on an
affine geometry, by which we mean the charts take values in $\R^n$ viewed as
an affine space and the transition maps are automorphisms of this affine space. The
fundamental object on such manifolds turns out to be a flat, torsion free
connection $\nabla$. This connection allows one to make sense of
differentiation on an affine manifold and hence make sense of the notion of a
parameterized curve having zero acceleration. Curves on an affine manifold
satisfying this property will be called \emph{geodesics}. Equivalently, a curve
on an affine manifold is a geodesic if it is a straight line in every affine
chart.

Of fundamental importance in the theory of affine manifolds is the case when
every geodesic can be continued indefinitely. Affine manifolds satisfying this
property are said to be \emph{complete}. In \cite{MarkusAuslander} Auslander
and Markus showed that completeness of an affine manifold $M$ is equivalent to
$M$ being the quotient of a discrete group of affine transformations acting
freely and properly on $\R^n$. There are many open questions surrounding
completeness of affine manifolds. For example, it was conjectured by Auslander
in \cite{AUSLANDER1964131} that the fundamental group of such an affine
manifold cannot be too "big" - namely that it should be solvable up to finite
index. If this conjecture holds, then it implies that a free group should not
be able to act on $\R^n$ by affine transformations such that the quotient is a
smooth, compact affine manifold.

The motivation behind this conjecture comes from the fact that affine geometry
shares many of the same properties as the more familiar Euclidean geometry. A
group which acts isometrically on $\R^n$ such that the quotient is a smooth,
compact manifold was shown to be abelian up to finite index by Bieberbach. He
classified these groups called \emph{Euclidean crystallographic groups} in
\cite{Bieberbach1910} and \cite{Bieberbach1912}. Unexpectedly, Margulis showed
in \cite{Margulis1983} that a free group could act freely and properly on
$\R^n$ by affine transformations, however, the quotient was not a
\emph{compact} manifold. Even though these affine manifolds failed to be
compact, they still garnered a lot of interest and were dubbed \emph{Margulis
    spacetimes}, owing to the fact that they inherit a unique parallel Lorentzian
metric.

Another open conjecture in affine geometry is the \emph{Markus conjecture}
which was posed by Markus in \cite{markus1963cosmological}. It states that if a
compact affine manifold possesses a volume form which is \emph{parallel} with
respect to the affine connection, then the affine manifold is complete. In
Chapter \ref{ch:nilpotent-hol}, we give the positive answer to this conjecture
in the case the fundamental group is nilpotent which was first proven in
\cite{nilpotent-holonomy}. This conjecture is open in full generality, however
it has been solved in some other special cases. For example, in
\cite{Carriere1989} they showed the conjecture holds in the case that the
affine manifold is a compact, flat Lorentzian manifold.

\section*{Integral affine manifolds and Lagrangian fibrations}

If an affine manifold (or its orientable double cover) admits a parallel volume
form, then the linear component of all its transition maps must live in the
group of invertible matrices of determinant $\pm 1$. An interesting subgroup of
this is the group $\mathrm{GL}(\Z^n)$, consisting of invertible \emph{integer}
valued matrices whose inverse is integer valued. Affine manifolds where the
linear component of their transition maps lie in this group are called
\emph{integral affine manifolds}. As well as having a characterization in terms
of admitting a flat, torsion-free connection whose holonomy lies in
$\mathrm{GL}(\Z^n)$, an integral affine structure on a manifold $M$ can also be
shown to be equivalent to a full-rank lattice in the cotangent bundle of $M$
that is locally spanned by closed one-forms. It is well known that the
cotangent bundle of any manifold admits a \emph{symplectic structure} and this
observation connects the study of integral affine manifolds to symplectic
geometry.

A \emph{symplectic manifold} is a manifold that admits a closed, non-degenerate
two-form called the symplectic form. By necessity, such a manifold is always
even dimensional. There is an important class of submanifolds of symplectic
manifolds called \emph{Lagrangian submanifolds}. These are submanifolds which
are half the dimension of the ambient manifold and such that the symplectic
form restricts to zero along them. If $B$ is any manifold, then its cotangent
bundle $T^*B$ is symplectic and the fibres of the projection map $\pi: T^*B \to
    B$ are Lagrangian submanifolds of the cotangent bundle. This leads into the
more general notion of a \emph{Lagrangian fibration}. If $M$ is a symplectic
manifold then a Lagrangian fibration is a surjection $\pi: M \to B$ where $B$
is any smooth manifold, such that the fibres of the projection are Lagrangian
submanifolds.

Of particular interest is when $\pi$ is a surjective submersion and the fibres
are compact and connected. In this case, the fibres are tori and the base space
$B$ inherits the structure of an integral affine manifold. The integral affine
structure on the base turns out to be an invariant of the Lagrangian fibration.
Moreover, the lattice in the cotangent bundle of $B$ which induces this
integral affine structure gives rise to a canonical Lagrangian torus fibration
over $B$. This fibration is a local model for all Lagrangian torus fibrations
over $B$ which induce the given integral affine structure. The obstruction to a
generic Lagrangian torus fibration being globally equivalent to this local
model was studied in detail by Duistermaat in \cite{ActionAngle}. He defined a
characteristic class called the \emph{Chern class} and its vanishing gives rise
to a global section of the Lagrangian fibration which can be used to identify
it with the local model.

This identification is only topological in the sense that the symplectic
structure of the fibration is not preserved. However, if the global section
determines a Lagrangian submanifold of the total space of the fibration then
the fibration identifies both topologically and symplectically with the local
model as shown also in \cite{ActionAngle}. This was expanded upon further by
Zung in \cite{zung2001symplectictopologyintegrablehamiltonian}. He defined
another characteristic class called the \emph{Lagrangian class} which maps onto
the Chern class through a certain group homomorphism. If this characteristic
class vanishes then the Lagrangian fibration will be both topologically and
symplectically equivalent to the local model. \newpage

\section*{Outline}

In this section we will briefly present the main structure of the thesis which
can hopefully be used as a guide to the reader. The primary goal of this thesis
is to classify the Lagrangian fibrations over closed surfaces. Our path to
achieving this classification will require us to use tools from many different
fields of mathematics such as symplectic geometry, algebraic topology and group
theory.

In Chapter \ref{ch:prelim} we develop some results from the general theory of
(integral) affine manifolds, the representation theory of groups and group
cohomology. We describe what is meant by \emph{completeness} of an affine
manifold and give two equivalent perspectives on this property. The first view
point introduces the notion of the \emph{developing map} which is a local
diffeomorphism from the universal cover of the affine manifold into $\R^n$. An
affine manifold is then called complete if this developing map is a global
diifeomorphism. The other equivalent definition of completeness is obtained by
noting that an affine manifold inherits a canonical flat, torsion-free
connection. Geodesic completeness of this connection is then equivalent to
completeness of the affine structure. We also introduce integral affine
manifolds which can be realised as a full-rank lattice in the cotangent bundle
of the manifold that is locally spanned by closed one-forms. These will become
key objects in the classification of Lagrangian fibrations as the base space
always inherits such a structure.

In Chapter \ref{LagrangianFibrations} we explore these objects in more detail.
It is well known that the cotangent bundle of any smooth manifold inherits the
structure of a \emph{symplectic manifold}. A symplectic manifold is a manifold
equipped with a closed, non-degenerate two form. The fact that an integral
affine structure is induced from a lattice in this cotangent bundle provides an
important link between integral affine geometry and symplectic geometry. The
main result of this section is that any Lagrangian fibration over a fixed base
$B$ gives $B$ the structure of an integral affine manifold. There is an action
of the cotangent bundle of $B$ on any Lagrangian fibration over $B$ and the
fixed-point set of this action corresponds to a lattice which induces the
integral affine structure.

In Chapter \ref{ch:nilpotent-hol} we start developing the tools required to
classify the integral affine structures on closed surfaces. The main result of
this section is that if the fundamental group of a compact affine manifold is
nilpotent, completeness of the affine structure is equivalent to the existence
of a volume form which is parallel with respect to the canonical affine
connection. This is a partial resolution to a conjecture of Markus in
\cite{markus1963cosmological} which states that this should hold true even when
the fundamental group is not nilpotent.

Chapter \ref{ch:class-IAS} contains the classification of the integral affine
structures over closed surfaces. This is one of the main invariants of a
Lagrangian fibration over these surfaces. Using the tools built up in Chapter
\ref{ch:prelim} and Chapter \ref{ch:nilpotent-hol} we are able to reduce the
problem of classifying these structures to a relatively simple computation
involving linear algebra and group theory. We also make use of a powerful
theorem proved by Milnor in \cite{Milnor1957/58} which states that a closed,
orientable affine surface must have genus less than $2$.

In the last chapter, we complete the classification of the Lagrangian
fibrations over the Klein bottle and the torus. In order to do this, we find
two more invariants of the fibration. These are the \emph{Chern class} and the
\emph{Lagrangian class}. The Chern class determines the topological nature
while the Lagrangian class is a refinement of this and determines the
symplectic structure of the fibration. These invariants are defined in terms of
sheaf cohomology which can make explicit calculation of them difficult in
practice. However, we show that these cohomology group can be computed using
group cohomology which we developed in Chapter \ref{ch:prelim}. We calculate
these explicitly for both the torus and the Klein bottle.

Finally, we show how to explicitly construct a Lagrangian fibration over the
torus or Klein bottle from their invariants. This is done using a symplectic
version of Dehn surgery called Luttinger surgery and first described in
\cite{LuttingerSurgery}. This completes the main content of this thesis and
gives an explicit way to construct any given Lagrangian torus fibration over a
closed surface.